\documentclass{beamer}
\usepackage{moreverb} 
\usepackage{listings}
\usepackage{mflogo}
% imprimir
% \documentclass[handout]{beamer} 
% \usepackage{pgfpages}
% \pgfpagesuselayout{4 on 1}[a4paper,landscape,border shrink=5mm]

\mode<presentation> {
  \usetheme{Warsaw}
  \setbeamercovered{transparent}
}

\usebackgroundtemplate{\includegraphics[width=\paperwidth]{format/libresoft-bg.png}}
% \usepackage[spanish]{babel}
\usepackage[utf8]{inputenc}
\usepackage{graphics}
\usepackage{amssymb} % Simbolos matematicos

%%%%%%%%%%%%%%%%%%%%%%%%%%%%%%%%%%%%%%%%%%%%%%%%%%%%%%%%%%%%%%%%%%%%%%%%%%%%%%%%%%%%%%%
%%
%% alert for highlighting
%% textit for Italic
%% texttt for URLs, etc
%%
%%%%%%%%%%%%%%%%%%%%%%%%%%%%%%%%%%%%%%%%%%%%%%%%%%%%%%%%%%%%%%%%%%%%%%%%%%%%%%%%%%%%%%%

%% Metadatos del PDF.
\hypersetup{  
  pdftitle={Project Evaluation Assignment},
  pdfauthor={Daniel Izquierdo, Pedro Coca},
  pdfcreator={GSyC/Libresoft},
  pdfproducer=PDFLaTeX,
  pdfsubject={Master on Free Software},
}
%%

\defbeamertemplate*{footline}{shadow theme}
{%
  \leavevmode%
  \hbox{\begin{beamercolorbox}[wd=.5\paperwidth,ht=2.5ex,dp=1.125ex,leftskip=.3cm plus1fil,rightskip=.3cm]{author in head/foot}%
    \usebeamerfont{author in head/foot}\insertframenumber\,/\,\inserttotalframenumber\hfill\includegraphics[scale=0.40]{format/cc-by-80x15.png} \hspace{0.1cm}\insertshortauthor 
% \usebeamerfont{author in head/foot} \includegraphics[width=0.7cm]{format/cc-by.png} \hfill\insertshortauthor
  \end{beamercolorbox}%
  \begin{beamercolorbox}[wd=.5\paperwidth,ht=2.5ex,dp=1.125ex,leftskip=.3cm,rightskip=.3cm plus1fil]{title in head/foot}%
    \usebeamerfont{title in head/foot}\insertshorttitle%
  \end{beamercolorbox}}%
  \vskip0pt%
}

\begin{document}

\title{Project evaluation assignment}
\subtitle{Master on Libre Software 2011-12}
\institute{\texttt{http://master.libresoft.es} \\ Identi.ca - Twitter: \texttt{@mswl}}
\author{Daniel Izquierdo, Pedro Coca} 
%\date{\today}
\date{November 11th, 2011}

\frame{
\maketitle
\begin{center}
\includegraphics[width=6cm]{format/gsyc-urjc}
\end{center}
}

%% License slide
\begin{frame}
  \vspace{2cm}
  \begin{center}
    {\small (cc) 2011 Daniel Izquierdo, Pedro Coca, LibreSoft} \\
%    \vspace{0.25cm}
    \medskip
    {\scriptsize This work is licensed under \\ a Creative Commons Attribution 3.0 License}
%    \vspace{0.10cm}
  \end{center}
  \begin{center}
    \href{http://creativecommons.org/licenses/by/3.0/es}{\includegraphics[width=2cm]{format/cc-by.png}} \\
    {\tiny \url{http://creativecommons.org/licenses/by/3.0}}
  \end{center}
\end{frame}%%

\usebackgroundtemplate{}

\AtBeginSubsection[]
{
  \begin{frame}<beamer>{Table of Contents}
    \tableofcontents[currentsection,currentsubsection]
  \end{frame}
}
p
%%%%%%%%%%%%%%%%%%%%%%%%%%%%%%%%%%%%%%%%%%%%%%%%%%%%%%%%%%%%%%%%%%%%%%%
\section{Assignment}
%%%%%%%%%%%%%%%%%%%%%%%%%%%%%%%%%%%%%%%%%%%%%%%%%%%%%%%%%%%%%%%%%%%%%%%

\subsection{Project evaluation}

\begin{frame}
\frametitle{Project evaluation}

\begin{itemize}
\item Analise several FLOSS tools 
\item Use a lightweight evaluation framework
\begin{itemize}
	\item OpenBRR
	\item QSOS
\end{itemize}

\item Source Code Management tools:
   \begin{itemize}
   \item Subversion
   \item Git
   \item Mercurial
   \end{itemize}

\end{itemize}
\end{frame}

%%%%%%%%%%%%%%%%%%%%%%%%%%%%%%%%%%%%%%%%%%%%%%%%%%%%%%%%%%%%%%%%%%%%%%%


\subsection{Role Play}

\begin{frame}
\frametitle{Role Play}

\begin{itemize}
\item Based on the evaluation performed with the SCM tools
\item Role: "Head of Development Department" 
\item Short report justifying what SCM will be the solution for your department
\item State any assumptions considered relevant for your decision: 
\begin{itemize}
	\item Programming language team expertise 
	\item Integration with other tools
	\item ...
\end{itemize}

\end{itemize}

\end{frame}

%%%%%%%%%%%%%%%%%%%%%%%%%%%%%%%%%%%%%%%%%%%%%%%%%%%%%%%%%%%%%%%%%%%%%%%
\section{References}
%%%%%%%%%%%%%%%%%%%%%%%%%%%%%%%%%%%%%%%%%%%%%%%%%%%%%%%%%%%%%%%%%%%%%%%

\begin{frame} 
\frametitle{References}

\begin{itemize}

\item SCM tools definition
        \begin{itemize}
	\item \tiny{\url{http://en.wikipedia.org/wiki/Source\_Code\_Control\_System}}
        \end{itemize}

\item Source Code Management tools websites:

        \begin{itemize}
        \item Subversion: \small{\url{http://subversion.apache.org/}}
        \item Mercurial: \small{\url{http://mercurial.selenic.com/}}
        \item Git: \small{\url{http://git-scm.com/}}
        \end{itemize}

\end{itemize}

\end{frame}

%%%%%%%%%%%%%%%%%%%%%%%%%%%%%%%%%%%%%%%%%%%%%%%%%%%%%%%%%%%%%%%%%%%%%%%

% Final slide
\usebackgroundtemplate{\includegraphics[width=\paperwidth]{format/libresoft-bg}}

\frame{
\maketitle
\begin{center}
\includegraphics[width=6cm]{format/gsyc-urjc}
\end{center}
}

\end{document}

